% Auto-generated by src/evaluation/analysis.py for ACL 2027 paper

% --- category_performance_table ---
\begin{table*}[t]
\centering
\small
\begin{tabular}{lrrrrrr}
\toprule
\textbf{Syntactic Category} & \textbf{Instances} & \textbf{Pairs} & \textbf{Weighted F1} & \textbf{SoftCons} & \textbf{HardCons} & \textbf{RCSI (\%)} \\
\midrule
Adverbial & 4 & 2 & 1.0000 & 1.0000 & 1.0000 & 0.0\% \\
Nominal & 7 & 3 & 0.8000 & 1.0000 & 1.0000 & 0.0\% \\
Prepositional & 6 & 3 & 1.0000 & 1.0000 & 1.0000 & 0.0\% \\
Quantificational & 6 & 3 & 0.5556 & 0.3333 & 0.3333 & 100.0\% \\
Verbal & 8 & 4 & 0.7500 & 0.5000 & 0.5000 & 66.7\% \\
\midrule
\textbf{Overall} & \textbf{31} & \textbf{15} & \textbf{0.8260} & \textbf{0.7333} & \textbf{0.7333} & -- \\
\bottomrule
\end{tabular}
\caption{Disaggregated performance across PhrasIS syntactic chunk categories for crl on track2.}
\label{tab:syntactic_breakdown}
\end{table*}

% --- error_cases_table ---
\begin{table*}[t]
\centering
\footnotesize
\begin{tabular}{p{2.0cm}p{3.2cm}p{3.2cm}ccp{5.0cm}}
\toprule
\textbf{Category} & \textbf{Phrase A} & \textbf{Phrase B} & \textbf{Gold (Dir/Rev)} & \textbf{Pred (Dir/Rev)} & \textbf{Diagnosis} \\
\midrule
Quantificational & nueve niños entusiastas & varios niños & FOR/BAC & FOR/FOR & Reversal Curse Entailment Loop in Quantificational Chunk: Monotonicity entailment from specific to general bound correct... \\
Quantificational & todos asistieron & alguien asistió & FOR/BAC & FOR/FOR & Reversal Curse Entailment Loop in Quantificational Chunk: Monotonicity entailment from specific to general bound correct... \\
Verbal & están bailando con alegría & se están moviendo & FOR/BAC & FOR/FOR & Reversal Curse Entailment Loop in Verbal Chunk: Model detects event subsumption ('están bailando con alegría' -> 'se est... \\
Verbal & caminando muy despacio & caminando despacio & FOR/BAC & BAC/BAC & Reversal Curse Entailment Loop in Verbal Chunk: Model detects event subsumption ('caminando muy despacio' -> 'caminando ... \\
\bottomrule
\end{tabular}
\caption{Representative qualitative failure cases and linguistic error diagnoses.}
\label{tab:qualitative_error_cases}
\end{table*}

